\subsubsection{A 类：先交代机制，再讨论数值}

A 类行文从研究对象、坐标、尺度和守恒关系进入。高信息密度表达包括“以……为
研究对象，取……方向为正”“由质量/能量/动量守恒可得”“在边界……处满足……”
“将区域离散为……，当网格由……加密至……时，目标量变化……”。推导篇幅不必刻意拉长，
假设、方程、初边值和数值格式仍要彼此闭合。结果解释需回到物理量、数量级和机制；
一条平滑曲线本身不足以说明模型正确。

需要避开的写法是先报“使用微分方程模型”，随后给软件输出；或者给出方程
却没有边界、单位、稳定条件和网格收敛。图形以结构示意、网格、状态曲线和残差/
收敛图为主，颜色只服务于变量区分。

A070、A147、A195、A212、A028、A115、A217、A001、A022、A171、A0127、A0165、
A0175、A092、A016、A053、A163、A178、A242 与 A196 共二十篇 A 题均已通过人工检查。
这些论文最值得学习的共同点，是把“为什么可以这样算”留在正文中：守恒关系授权方程，
边界条件限定数值格式，几何次序筛去无意义的根，现实动作排除数学上可行但物理上不可能的
分支。作者不急于给算法名，而是先写观察、冲突或单调关系，再自然地落到求解方法。

A028 与 A115 的可迁移价值在于把理想几何、工程约束和执行方案分层表达。前者先以坐标和
面板关系构造几何候选，再用索长、邻距和受力兼容性判断能否执行；后者把节点逼近、连续
工作面和接收效率分成三个统计定义，提醒写作者不要用一个代理量代替完整物理链。仿写时应保留
“先得到候选，再检查……”这一顺序，并在结果段同时交代几何读数、约束余量和工程含义。

A053、A163、A178 与 A242 围绕同一类板凳龙任务，却给出了四种很像人类研究的推进方式。
A053 用完整反证把首次碰撞逐节前推到龙头；A163 先证明自由度，再由全局轨迹图缩小局部
搜索域；A178 让代理圆只负责夹住区间，最终仍回到四角点精判；A242 先给物理下界和经验
可行上界，再在有来由的区间内优化。四篇都把“缩域依据”写在“搜索方法”之前，避免读者
看到一个凭空出现的上下界。

其中 A053 的结果解释尤其值得直接转化为写作习惯：参数只变化 $0.01\,\mathrm{cm}$ 时，先查
首次触发事件是否从一个构件切换到另一个构件，再解释输出为何跃变；某个终点方案看似可行，
则把它放回全过程寻找提前碰撞，以反例裁决简化方案；半径继续减小时若要求后续构件倒退，
就由这一现实后果推出数学下界。A196 则先证明目标区间连通，才使用二分；理论误差、输出
精度和多初值结果分开报告。它们共同说明，模型可信度常来自一两句关键判断，而不是更长的
算法介绍。

A217 展示了机理几何题更值得迁移的一面：先由轴对称抛物面的截面关系把焦距问题
压缩成一维驻点，再沿用同一坐标框架处理面板姿态和遮挡，最后以三角面板射线追迹
配合蒙特卡洛积分估计接收效率。正文不靠长串套话衔接，而用“注意到……，故……”
交代关键观察，用“考虑到……，本文将……”给出取舍；公式之后马上说明几何量的
含义，成对图形则以“图左/图右”分别解释基准状态与调整状态。这种写法把观察、
选择、推导和图证压在同一条论证线上，比先报算法名称再补理由更像真实建模过程。

A001 展示了机理题中“公式后留一句人话”的用法。作者在 RK4 推进式后逐项说明起点、
中点和终点斜率，再用“易知，当上述斜率加权平均时，时间中点斜率权值最大”归结。
先观察前期位移不稳定，才截取 $100$--$180\,\mathrm{s}$ 计算平均功率；阻尼优化先用
增大步长粗略遍历，再用遗传算法求精。旋转阻尼影响小时，只作“可能是纵摇角速度较小”
的局部解释；周期误差和 $\pm1\%$ 转动惯量扰动分别承担检验与可靠性任务。这种
“现象—时间窗口—求解分工—检验规则”比写“使用 RK4 和 GA”更像人类作者推进。

A022 则把“边界先行”写得清楚：先计算平衡位高 $h_0$ 排除浸没和锥面暴露，再进入运动方程。第二问不是把复数解和数值积分并列，而是由前一问的数值结果授权改用：恒定阻尼可写出解，幂律阻尼则用 \texttt{trapz} 配合粗搜、局部加密和排序。问题三在相对垂荡很小时才把振子放入以浮子为非惯性系的纵摇算式，问题四再由图形发现旋转阻尼单调增长后设为上界并只局部加密纵摇阻尼。“观察—授权—改用—近似—边界”是比“列出 ode45、trapz、扫描”更像人类论文的连接。

A171 把机理题的“人话”写在方法切换之前：先算平衡压缩量和吃水，说明分段浮力在哪些区间适用；常阻尼是常系数方程，故列出拉普拉斯变换的三步，幂律阻尼解析解难求，才改用四阶 Runge--Kutta。优化段先定义瞬时功率、最后十个稳定周期和约束，再给阻抗匹配或梯形积分的结果。每种情形都用 $10,20,40,60,100\,\mathrm{s}$ 交付同一组状态量，读者可以逐项对照；检验段又把初始过渡误差和稳态误差分开，最后落到“可以接受”的判断。可迁移的不是“任何题都用 RK4”，而是“边界—方程性质—方法—同口径结果—判定”的段落顺序。

A0127 把工程计算中容易被 AI 写空的“为什么这样算”保留下来。问题分析先写“它实际上是一道优化问题”，随后将第二问称为第一问的反问题、第三问写成放宽约束；读完题面公式后，又逐项判断哪些量可直接获得，最后才说计算关键是五类光学效率。塔影没有用“为简化计算”一笔带过，而是先说明每面镜姿态不同带来的边界计算负担，再以完全覆盖镜数量占优、镜面中心对称和边界附近统计补偿，推出镜面中心是否落入阴影区的二值规则。后面的空间云图也不是装饰：北侧效率较高，于是镜面向北集中、塔向南移动，塔位变量随之限定在负 $Y$ 轴。搜索段同样先写直接高精度遍历的计算量，再让二分、大步长和小步长分别承担定位区间与提高精度。可迁移的是“问题身份—剩余未知—简化依据—可执行判定—空间证据—计算预算”的自然推进，不是机械复用“考虑到”“因此”等词。

A0165 又展示了另一种更像人工建模记录的推进。摘要算法名不应先于把定日镜计算拆成太阳位置、镜面法向、遮挡、截断和功率五个传播阶段，分别用“确定”“计算”“判断”和“从而得到”交代输入输出关系。几何准备段先说“可将矩形分为两个三角形”，随后将点在矩形内的问题等价成两次三角形判定；遮挡段只检查六个近邻，并用 $2\times2$ 栅格构造快速计算核，且报告其与完整计算的偏差小于 $2\%$。优化段分析一个变量时明确固定其余变量，由单峰曲线才采用三分搜索；空间散点图显示北侧效率占优后，塔位设定依据应来自向南移动。第三问沿用第二问时又把“保留什么”和“新求什么”分开：镜宽、塔位和布场被冻结，只搜索各圈高度。检验也不能只报一个误差数，还要先预言峰值应随太阳视运动移动，主动改换五个采样时刻，再比较散点峰值的迁移方向。值得迁移的是“过程分段—等价转化—快速核及偏差—固定变量—图形触发设计—继承与冻结—结构预言检验”的信息顺序。

A0175 更像一份把程序控制流翻译成论文的人工记录。塔影段先写“若塔身对某入射光线造成阴影遮挡损失，那么不再考虑……”；只有塔影未发生，才继续判断入射阴影、反射遮挡和集热器接收。邻接矩阵也不是突然出现：作者先将 1745 面镜、镜面网格和锥束角度共同造成的计算量写明，再用遮挡主要发生在邻近镜之间这一物理理由，把候选限制在 20 m 内。步长由 6 m 逐步降到 0.25 m，单位面积输出仍在上升，正文只说从 1.5 m 起“增长的速度变得缓慢”，因而“认为……接近稳定”；这种有限语气比“充分收敛”更符合证据。问题二先以 $10\%$ 镜面重复 10 次估计粗参数，遍历旋转角和高度发现影响不大后冻结，随后将预算转向镜宽和塔位，最后用 $30\%$ 镜面重复 5 次汇总。表 7 得到 $68.2443\,\mathrm{MW}$ 后，作者没有停在报数，而是从“额定功率只需要达到即可”继续推理：删去高遮挡镜面会同时降低总面积、总功率和相互遮挡，因此应进一步兼顾单位面积输出。值得迁移的是“条件—停止—继续”“计算量—局部性—邻接”“序列—趋缓—工作步长”“弱影响—冻结—转移预算”和“达标结果—连锁作用—目标修正”。

A092 的自然感来自作者愿意把“还剩什么没算出来”写清楚。把题面数据代入第一问模型后，大气透射、余弦与遮挡等量可以直接获得，正文才说“唯一无法直接求得的是截断效率”，并用随机入射方向、随机镜面点和 10,000 次接收判定解释蒙特卡洛的职责。入射遮挡与反射遮挡可能重叠时，作者也没有直接相加面积，而是先指出重复计数风险，再离散镜面、统计两区域并集内的点。结果图的推进同样分两步：环序号曲线只负责发现北侧高值并写“由此推测”，随后专门选春分日上午 10 时画二维分布，才把空间猜想转成布局证据。第二问固定其余参数，分别沿东西和南北扰动塔位；东西响应关于原点近似对称，上午和下午的影响又相互抵消，因此横向坐标被主动冻结。第三问则从前问的同心圆结果继续推理：同心圆各向同性，同一圈可共享镜面尺寸和高度，逐镜变量便压成逐圈变量，求解方法也由高维 GA 切换为变步长遍历。这里最值得写入 Skill 的不是算法名称，而是“逐项排除—唯一剩余量—随机估计”“曲线—空间猜想—代表截面复核”“方向扰动—对称抵消—冻结变量”和“前问结构—对称依据—按组降维”的段落动作。论文将方案与 100 组随机可行解比较时，结论只应保留在该样本背景内。

A016 的人类感来自作者把数学多解、计算缩域和数值回退都写成了可见判断。相邻把手定长方程有多根时，正文先说板凳依次排列，故角度必须满足前后顺序；又因相邻把手只转过一个小角度，才在可行根中取最近根。碰撞段也不是泛称“降低计算复杂度”，而是先用下游板凳沿前件轨迹运动的事实，判断最早碰撞主要涉及龙头或第二板凳，再列出真正要检查的顶点；可能晚些出现的特殊碰撞，则由“按照时间遍历，找到开始碰撞的时刻立即停止”排除。

逐步逼近和首次碰撞搜索都保留“试一步—越界—退一步—缩步—再判断”，最小步长随后被译成误差小于其两倍。面对题目要求的“调头曲线还能否缩短”，作者没有为了交付一个优化数值而强行搜索，而是利用相切和比例关系证明两段圆弧总长度固定；证明没有自由度后，才继续计算题设比例下的调头路径。问题五同样先证明任一把手速度变化 $k$ 倍会沿链条传递，再整体缩放问题四的速度矩阵。这里值得迁移的是“物理顺序—选根”“最早事件—缩域”“越界—回退—误差”“固定量—不可优化”和“局部齐次—整体缩放”，而并非单独模仿“因此、可以发现、可见”。

